\documentclass[10pt,a4paper]{amsart}
\usepackage[margin=19mm]{geometry}
\usepackage{amsmath,amssymb,mathtools}
\usepackage[scr=rsfs,cal=euler]{mathalfa}
\usepackage{xfrac}
\usepackage[unicode,hidelinks,bookmarks=false]{hyperref}
\newcommand{\ie}{i.e.\ }
\newcommand{\cf}{cf.\ }
\newcommand{\resp}{respectively\ }
\newcommand{\Subsection}{\subsection}
\providecommand{\nobreakdash}{\nobreak\hbox{-}}
\setlength{\emergencystretch}{2em}
\multlinegap0pt
\usepackage{siunitx}
\usepackage{siunitx}
\usepackage{siunitx}
\usepackage{siunitx}
\usepackage{amssymb}
\usepackage{siunitx}
\newtheorem{theorem}{Theorem}
\title{Geometric Factorization of Sufficient Harmonic Representations}
\author{Kennon Stewart}
\date{Source: arXiv:2606.07346v1 (CC BY 4.0)}
\begin{document}
\maketitle

\noindent\emph{Source and licence.} Geometric Factorization of Sufficient Harmonic Representations. Version arXiv:2606.07346v1 (CC BY 4.0), distributed by arXiv under the Creative Commons Attribution 4.0 International licence, http://creativecommons.org/licenses/by/4.0/, as recorded in the arXiv licence metadata for this exact version and shown on the abstract page https://arxiv.org/abs/2606.07346v1. Full licence notice: \texttt{licenses/arxiv-2606.07346-CC-BY-4.0.txt}.\par\smallskip

\noindent\textit{Editorial scope.} Counted unit: the article's theorem theorem (source0.tex lines 237--247). The article's complete original proof of it is delivered with it (source0.tex lines 249--373, one \texttt{proof} environment of 125 lines). No mathematics is written, repaired or re-derived: the statement and the proof are the article's own lines, in the article's own notation, and no retained line is substituted or re-flowed.  No further source-local context is needed: every cross-reference of every retained line resolves inside the two delivered spans.  Nothing was omitted from any retained span, so the package carries no omission record.  No retained line contains a citation, so the wrapper carries no bibliography and no bibliography entry.  No retained line contains a figure environment, an image-inclusion command or drawing code, so no image is delivered and the package declares no asset dependency.  Editorial formatting only, in addition: an AMS \texttt{amsart} preamble that restates the article's own declarations and macros used by the retained lines, namely the environments \texttt{theorem} on separate counters, together with the standard packages those lines need.  The retained environments print in this document as Theorem 1 in order of appearance, whereas the article prints the same items with the numbering of its own full document.  The complete native source of the article as distributed in the arXiv e-print for this exact version is bundled unchanged as \texttt{sources/202251/source0.tex}.
\begin{theorem}[Minimal Sufficiency of Harmonic Statistics]
Suppose the finite-band harmonic exponential family above is full rank: after choosing real coordinates for the matrix coefficients, the natural statistic has no nontrivial affine dependence almost surely, and the natural parameter space contains an open set. Then, for i.i.d. observations $X_1,\dots,X_n\in G$, the statistic
$$
T_N(X_1,\dots,X_n)
=
\left(
\sum{i=1}^n \rho_\lambda(X_i)^\dagger
\right){\lambda\in\widehat G_N}
$$
is minimal sufficient for $\theta$.
\end{theorem}

\begin{proof}
For i.i.d. observations \(X_1,\dots,X_n\in G\), the joint likelihood is
\[
L(\theta\mid X_1,\dots,X_n)
=
\prod_{i=1}^n p_\theta(X_i)
=
\exp\left[
\sum_{i=1}^n E_\theta(X_i)-nA(\theta)
\right].
\]
Substituting the finite-band harmonic energy gives
\[
\sum_{i=1}^n E_\theta(X_i)
=
\sum_{i=1}^n
\sum_{\lambda\in\widehat G_N}
d_\lambda
\operatorname{Re}
\operatorname{Tr}
\left(
C_\lambda^\dagger \rho_\lambda(X_i)
\right).
\]
Since the index set \(\widehat G_N\) is finite, we may interchange the sums and use
linearity of the trace:
\[
\sum_{i=1}^n E_\theta(X_i)
=
\sum_{\lambda\in\widehat G_N}
d_\lambda
\operatorname{Re}
\operatorname{Tr}
\left(
C_\lambda^\dagger
\sum_{i=1}^n \rho_\lambda(X_i)
\right).
\]
Therefore
\[
L(\theta\mid X_1,\dots,X_n)
=
\exp\left[
\sum_{\lambda\in\widehat G_N}
d_\lambda
\operatorname{Re}
\operatorname{Tr}
\left(
C_\lambda^\dagger T_{N,\lambda}(X)
\right)
-
nA(\theta)
\right],
\]
where
\[
T_{N,\lambda}(X)=\sum_{i=1}^n \rho_\lambda(X_i).
\]
Thus the likelihood depends on the sample only through
\[
T_N(X)=
\left(
T_{N,\lambda}(X)
\right)_{\lambda\in\widehat G_N},
\]
and the Fisher-Neyman factorization theorem implies that \(T_N\) is sufficient.

It remains to show minimality. Let \(X=(X_1,\dots,X_n)\) and
\(Y=(Y_1,\dots,Y_n)\). The likelihood ratio is
\[
\frac{L(\theta\mid X)}{L(\theta\mid Y)}
=
\exp\left[
\sum_{\lambda\in\widehat G_N}
d_\lambda
\operatorname{Re}
\operatorname{Tr}
\left(
C_\lambda^\dagger
\left[
T_{N,\lambda}(X)-T_{N,\lambda}(Y)
\right]
\right)
\right],
\]
since the terms \(nA(\theta)\) cancel.

If \(T_N(X)=T_N(Y)\), then the exponent is zero and the likelihood ratio is
equal to \(1\), hence constant in \(\theta\).

Conversely, suppose the likelihood ratio is constant in \(\theta\). Then
\[
\sum_{\lambda\in\widehat G_N}
d_\lambda
\operatorname{Re}
\operatorname{Tr}
\left(
C_\lambda^\dagger
\left[
T_{N,\lambda}(X)-T_{N,\lambda}(Y)
\right]
\right)
\]
is constant as a function of the natural parameters
\((C_\lambda)_{\lambda\in\widehat G_N}\). Because the natural parameter space
contains an open set and the family is full rank after choosing real
coordinates, this is possible only if
\[
T_{N,\lambda}(X)-T_{N,\lambda}(Y)=0
\quad
\text{for all }
\lambda\in\widehat G_N.
\]
Hence \(T_N(X)=T_N(Y)\).

Therefore
\[
T_N(X)=T_N(Y)
\quad\Longleftrightarrow\quad
\frac{L(\theta\mid X)}{L(\theta\mid Y)}
\text{ is constant in }\theta.
\]
By the likelihood-ratio characterization of minimal sufficiency,
\(T_N\) is minimal sufficient.
\end{proof}

\end{document}
